%! Author = murfe
%! Date = 06.09.2026
\chapter{Evaluation}
\label{ch:evaluation}
Dieses Kapitel untersucht im Rahmen von \ref{item:tz4}, ob das implementierte Job-Processing-System die in \cref{ch:anforderungen} definierten
Zusicherungen auch unter den jeweils betrachteten Fehlerszenarien einhält.
Hierzu werden zunächst der konkrete Versuchsaufbau und anschließend die für die einzelnen Zusicherungen durchgeführten
Fehlerszenarien und deren Prüfkriterien dargestellt.
Darauf aufbauend wird die konkrete Umsetzung der Tests inklusive Testparameter beschrieben und die beobachteten Ergebnisse ausgewertet.
Den Abschluss bildet die Bewertung der formulierten Hypothesen.

\section{Versuchsaufbau}
\label{sec:versuchsaufbau}
Sofern für die einzelnen Tests nicht anders angegeben, wird die Anwendung für die Tests mittels \texttt{@SpringBootTest}
vollständig mit allen ihren Komponenten gestartet.
Die Datenbank läuft dabei lokal über Docker auf einem Postgres 17 image.
Vor den jeweiligen Testdurchläufen wird sichergestellt, dass die relevanten Datenbanktabellen leer sind, sodass
die einzelnen Testläufe unabhängig voneinander sind und sich nicht gegenseitig beeinflussen.

Die verwendeten Systemparameter sind für die getesteten Szenarien angepasst.
Das Ziel dabei ist, die einzelnen Komponenten so aufeinander abzustimmen, dass die betrachteten Fehlerinjektionen reproduzierbar
durchgeführt und die vorgesehenen Konkurrenzsituationen gezielt erzeugt werden können.
Konkret wurde die Anzahl der Worker auf 1 gesetzt, die Anzahl der maximalen Verarbeitungsversuche für Jobs auf 5, die Lease-Dauer
auf 10 Sekunden, die Verarbeitungszeit der Jobs auf 5 Sekunden und das Recovery-Intervall ebenfalls auf 5 Sekunden.
Das Retry- und Backoff-Verhalten ist mit 200ms und einem Faktor von $2^n$ gewählt, um die Laufzeit der Tests zu begrenzen.
Findet ein Worker keinen freien Job, pausiert er und wartet 2 Sekunden, bis er den Job-Bestand erneut prüft.
Einzelne Parameter können für die unterschiedlichen Tests angepasst werden, so beispielsweise die Anzahl der Worker für die
nebenläufigkeitsabhängigen Tests.
Diese Änderungen werden bei der jeweiligen Testbeschreibung dargelegt.

\section{Fehlerszenarien und Operationalisierung}
\label{sec:fehlerszenarien}
Ausgangspunkt der folgenden Testszenarien sind die bereits in \cref{ch:anforderungen} definierten Zusicherungen.
Für jede Zusicherung wurde untersucht, welches konkret beobachtbare Systemverhalten ihre Verletzung anzeigen würde.
Die resultierenden Prüfkriterien und Fehlerinjektionen sind in \cref{tab:testmatrix} zusammengefasst.
Ziel dabei war, die Szenarien so zu entwerfen, dass möglichst realistische Ausfallbedingungen simuliert werden, wie
beispielsweise ein echter Netzwerkfehler oder ein Ausfall einer einzelnen Komponente.
Für jede Zusicherung wird mindestens ein Szenario durchgeführt, anhand dessen die jeweils relevanten Prüfkriterien
ausgewertet werden können.
Die Auswahl inklusive Begründung der betrachteten Fehlerklassen wurde bereits in \cref{sec:abgrenzungFehlermodell} beschrieben.

%! suppress = EscapeAmpersand
\begin{longtable}{@{}p{0.045\textwidth}p{0.20\textwidth}p{0.33\textwidth}p{0.34\textwidth}@{}}
    \toprule
    \textbf{Z} & \textbf{Fehlerinjektion} & \textbf{Umsetzung} & \textbf{Prüfkriterium} \\
    \midrule
    \endfirsthead

    \multicolumn{4}{c}{-- Fortsetzung von vorheriger Seite --} \\
    \toprule
    \textbf{Z} & \textbf{Fehlerinjektion} & \textbf{Umsetzung} & \textbf{Prüfkriterium / Metrik} \\
    \midrule
    \endhead

    \midrule
    \multicolumn{4}{r}{\footnotesize Fortsetzung auf nächster Seite} \\
    \endfoot

    \bottomrule
    \addlinespace
    \caption[Testmatrix der Evaluation]{Testmatrix der Evaluation mit Zuordnung der Systemzusicherungen zu den untersuchten
    Fehlerbedingungen, den eingesetzten Fehlerinjektionsmethoden und den jeweiligen Prüfkriterien.}
    \label{tab:testmatrix} \\
    \endlastfoot

    \ref{zus:atomare_persistenz} & Ungültige Payload & Fachlich und syntaktisch invalide Requests senden & \texttt{HTTP 400}, kein DB-Eintrag; Zeilenzahl vor/nach unverändert \\
    \addlinespace
    \ref{zus:duplikaterkennung} & Acknowledgement geht verloren & Abfangen und verwerfen des Acknowledgements, stattdessen Fehler werfen & Client startet automatischen Retry. Die \gls{api} wird insgesamt 2 Mal angesprochen. Es wird genau ein Job im Backend gespeichert und der Statuscode der vom Client erhaltenen Antwort ist \texttt{200 OK} \\
    \addlinespace
    \ref{zus:exklusive_rechte} & (a) Viele Worker konkurrieren um einen Job; (b) Zwei Worker versuchen gleichzeitig Ergebnis zu speichern & (a) $K$ Worker-Instanzen versuchen den gleichen Job zu beanspruchen, loggen des jeweils beanspruchten Jobs; (b) Synchronisation der Worker vor anschließendem parallelen Speicherversuch & (a) Von $K$ gleichzeitigen Versuchen darf genau einer erfolgreich sein; (b) Es gibt genau einen Speichervorgang zu dem Job. Verspäteter Aufruf wird per Fencing abgelehnt \\
    \addlinespace
    \ref{zus:atomare_wirkung} & Erzwungener Fehler zwischen Save und Commit & Test-Hook wirft Exception nach \textit{setStatus()}, vor \mbox{\textit{resultRepository.save()}} & Nach Fehler: kein Ergebnis $\wedge$ \textit{Running} \\
    \addlinespace
    \ref{zus:endzustaende} & Recovery-Scan parallel zu \mbox{\textit{finishJob}} & Recovery-Scan und \textit{finishJob} zum gleichen Job über Barrier synchronisieren und parallel durch verschiedene Threads ausführen & Kein Übergang aus \mbox{\textit{Succeeded}}/\textit{Failed} zurück in \textit{Queued} \\
    \addlinespace
    \ref{zus:retries} & (a) Transiente DB-Störung; (b) permanenter Fachfehler & (a) DB-Verbindung kurzzeitig kappen; (b) Anfrage mit existierendem Idempotenzschlüssel als permanenter Fehler & (a) Erfolg nach Wiederherstellung, Retry-Anzahl im Log; (b) sofortiger Abbruch ohne Retry-Schleife \\
    \addlinespace
    \ref{zus:verarbeitungsgarantie}, \ref{zus:idempotenz}, \ref{zus:terminierung} & Harter Worker-Ausfall während \textit{Running} & Worker-Instanz während Bearbeitung hart beenden & Job wird nach Recovery von einem neuen Worker aufgenommen und erreicht \mbox{\textit{Succeeded}}, gespeichertes Ergebnis entspricht dem erfolgreichen (zweiten) Ausführungsversuch \\
    \addlinespace
    \ref{zus:idempotenz} & (a) Recovery eines Jobs im Status \textit{Running}; (b) Lease abgelaufen ohne vorherige Recovery & (a) Manuelle Recovery eines aktiven Jobs. Speicherversuch durch nicht-besitzenden Worker vor und nach Abschluss durch besitzenden Worker; (b) \texttt{lease\_until} auf Vergangenheit setzen & (a) Zugriff wird abgelehnt, es wird genau ein Ergebnis vom besitzenden Worker gespeichert; (b) Zugriff wird abgelehnt, kein Ergebnis wird gespeichert \\
    \addlinespace
    \ref{zus:terminierung} & Job schlägt dauerhaft fachlich fehl & Processor-Mock wirft für bestimmte Job-ID bei jedem Versuch eine Exception & \texttt{attempt\_count} erreicht Maximum, Job geht in \textit{Failed} über; Zeit bis \textit{Failed} $\leq$ \texttt{maxAttempts} $\times$ (Lease + Recovery-Intervall + Worker idle-Timeout) \\
\end{longtable}

Eine Zuordnung der in \cref{tab:testmatrix} beschriebenen Fehlerszenarien zu den konkret implementierten Testklassen und
-methoden befindet sich im Anhang in \cref{tab:zuordnung_zusicherungen_tests}.

\section{Durchführung}
\label{sec:durchfuehrung}
Dieser Abschnitt beschreibt, wie die in \cref{tab:testmatrix} definierten Fehlerinjektionen konkret technisch umgesetzt wurden.

\subsection{Joberstellung und Idempotenz (\ref{zus:atomare_persistenz}, \ref{zus:duplikaterkennung})}
\label{subsec:durchfuehrung_z1_z2}
Die Überprüfung von \ref{zus:atomare_persistenz} erfolgt, indem eine syntaktisch sowie eine fachlich ungültige Anfrage über den Client an das laufende
Backend gesendet werden.
Vor und nach beiden Anfragen wird die Anzahl der im Job-Store vorhandenen Jobs abgefragt und verglichen.

Für \ref{zus:duplikaterkennung} wird das Verlorengehen einer Erfolgsantwort simuliert, damit der Client automatisch einen
Retry mit demselben Job durchführt.
Hierzu wird am \texttt{BackendRestClient} ein \texttt{ClientHttpRequestInterceptor} registriert, der beim ersten Aufruf die tatsächliche
Serverantwort abwartet, anschließend jedoch eine \texttt{IOException} auslöst, bevor die Antwort an den Client weitergereicht wird.
Dadurch verarbeitet das Backend die Anfrage und persistiert den Job, ohne dass diese Information beim Client ankommt.
Der clientseitige \texttt{RetryExecutor} sollte die Exception folglich als transienten Kommunikationsfehler interpretieren
und dieselbe Anfrage samt identischem Idempotenzschlüssel erneut übertragen.
Schließlich wird geprüft, ob tatsächlich nur ein Job persistiert wurde.
Auf diese Weise wird das in \cref{subsec:atLeastOnce} beschriebene Szenario eines verlorenen Acknowledgements durch
einen transienten Netzwerkfehler realitätsnah nachgebildet.
Zusätzlich prüft dieser Test die Durchführung eines automatischen clientseitigen Retries bei transienten Fehlern (vgl. \cref{sec:ergebnisse}).

\subsection{Retry-Verhalten bei Infrastrukturfehlern (\ref{zus:retries})}
\label{subsec:durchfuehrung_z6}
Für die Untersuchung transienter Datenbankfehler wird die über Docker bereitgestellte PostgreSQL-Datenbank in einem Testcontainer ausgeführt.
In der Literatur wird Toxiproxy als Werkzeug zur Fehlerinjektion bei Microservices empfohlen \parencite{miranda_proposed_2024}.
Daher fällt die Wahl auf Toxiproxy, um die Datenbankverbindung der Anwendung zu verwalten.
Die Fehlerinjektion erfolgt durch das gezielte Unterbrechen der über den Proxy geführten Verbindung, während der
Datenbank-Container selbst weiterläuft.
Ein Neustart der Datenbank selbst hätte bezogen auf den Test der Retry-Funktionalität gegenüber der gewählten Fehlerinjektion keinen
wesentlichen Mehrwert, da sowohl ein Neustart als auch ein hier simulierter Verlust der Datenbankverbindung aus Sicht der
Anwendung zunächst in einer nicht verfügbaren Datenbank resultieren.
Auch für die Dauerhaftigkeit und Konsistenz bereits erfolgreich persistierter Daten hätte ein Testfall, der den Crash oder
Neustart der Datenbank simuliert, keinen Mehrwert, da diese Eigenschaften nicht durch die entwickelte Architektur, sondern
durch die \gls{acid}-Eigenschaften von PostgreSQL sichergestellt werden.
Daher wird im Rahmen der Evaluation auf einen solchen Test verzichtet.
Da die Worker und der Recovery-Service für die folgenden Tests nicht benötigt werden, aber die Logs verunreinigen würden,
werden beide Komponenten für die in diesem Kapitel beschriebenen Tests deaktiviert.

Da eine bereits offene, im Verbindungspool vorgehaltene Datenbankverbindung von der Unterbrechung des Proxys nicht
zwangsläufig betroffen ist, wird der Verbindungspool nach der Unterbrechung zusätzlich gezielt geleert, sodass der nächste
Datenbankzugriff tatsächlich einen neuen, über den unterbrochenen Proxy geführten Verbindungsversuch erzwingt.
In Testfall (a) wird die Verbindung für eine festgelegte Dauer unterbrochen und anschließend automatisiert wiederhergestellt.
Als Variante hierzu wird zusätzlich untersucht, wie sich das System verhält, wenn die Unterbrechung dauerhaft bestehen
bleibt und sämtliche Wiederholungsversuche erschöpft werden, bevor die Datenbankverbindung wiederhergestellt wird.
Für Testfall (b) wird eine Anfrage mit einem bereits verwendeten Idempotenzschlüssel gesendet, um stellvertretend
die Behandlung eines nicht retrybaren Datenbankfehlers zu prüfen.

Die Retry-Logik wird dabei ausschließlich anhand der Joberstellung überprüft.
Sowohl das Beanspruchen als auch der Abschluss eines Jobs verwenden dieselbe \texttt{RetryExecutor}/ \texttt{DBRetryPolicy}-Komponente.
Eine gesonderte Untersuchung an diesen Stellen liefert daher keinen zusätzlichen Erkenntnisgewinn.
Der für \ref{zus:duplikaterkennung} durchgeführte Test belegt zudem das clientseitige Retry-Verhalten bei Kommunikationsfehlern zwischen Client und Backend.
Eine vollständige, zu (a) und (b) symmetrische Untersuchung der clientseitigen \texttt{RestMessageRetryPolicy} wird nicht zusätzlich durchgeführt,
da der Unterschied zur \texttt{DBRetryPolicy} lediglich in der konkreten Fehlerklassifikation liegt.

\subsection{Atomarität der Ergebnisspeicherung (\ref{zus:atomare_wirkung})}
\label{subsec:durchfuehrung_z4}
Um einen Fehler gezielt zwischen der Statusänderung und der Persistierung des Ergebnisses auszulösen, wird ein
Test-Hook in die Implementierung von \textit{finishJob} eingefügt, der im produktiven Betrieb keine Wirkung entfaltet und
im Testfall über Mockito so konfiguriert wird, dass er nach der Statusänderung, aber vor dem Speichern des Ergebnisses
eine Exception auslöst.
Anschließend wird geprüft, dass die Änderungen verworfen wurden und kein Zwischenzustand im Job-Store verblieben ist.

\subsection{Nebenläufigkeit beim Claiming und Abschluss (\ref{zus:exklusive_rechte}, \ref{zus:endzustaende})}
\label{subsec:nebenlaeufigkeits_fehlersimulation}
Die Testszenarien zu \ref{zus:exklusive_rechte} und \ref{zus:endzustaende} hängen vom zeitlichen Zusammenspiel mehrerer konkurrierender Operationen ab.
Um die konkurrierenden Aufrufe möglichst zeitgleich zu starten, werden die beteiligten Threads über eine \texttt{CyclicBarrier}
synchronisiert, sodass keiner der Aufrufe beginnt, bevor nicht alle beteiligten Threads bereitstehen.
Um die Worker synchronisieren zu können, wird im Rahmen der Tests ein separater Threadpool genutzt, dem die jeweiligen
Operationen übergeben werden.
Damit die Test-Jobs garantiert von den Operationen des Threadpools verarbeitet werden, wird die Anzahl der durch das Backend
gestarteten Worker in diesen Tests auf 0 gesetzt.

Da die getesteten Eigenschaften hier Zugriffs- und Schreibrechte der Worker auf spezifische Datensätze betreffen, ist
eine notwendige Voraussetzung für die Aussagekraft dieser Tests, dass die konkurrierenden Zugriffe tatsächlich auf
Datenbankebene aufeinandertreffen und nicht bereits vorher durch den Verbindungspool serialisiert werden.
Die verwendete maximale Poolgröße muss daher mindestens der Anzahl gleichzeitig zugreifender Worker entsprechen.
Um diese Systemeigenschaft sicherzustellen, wird dies vor Ausführung der betroffenen Testklasse automatisiert überprüft
und die Testausführung andernfalls abgebrochen.

Für \ref{zus:exklusive_rechte} (a) wird ein einzelner Job von mehreren Worker-Threads gleichzeitig beansprucht.
Anschließend wird geprüft, wie viele Worker dabei Erfolg hatten.
Da für den Test genau ein Job im Job-Store persistiert ist, sollte nur ein Worker erfolgreich sein und alle anderen
ein leeres \texttt{Optional}-Objekt erhalten.

\ref{zus:exklusive_rechte} (b) testet nach dem gleichen Prinzip das gleichzeitige Speichern eines Ergebnisses desselben Jobs.
Ein besitzender und ein nicht besitzender Worker versuchen über eine Barrier synchronisiert ein Ergebnis zu speichern.

Zum Test von \ref{zus:endzustaende} werden \textit{finishJob} und der direkte Aufruf von \textit{recoverJob} auf denselben
Job angewendet, wobei der Job absichtlich mit bereits erreichter maximaler Versuchsanzahl präpariert wird, sodass je nach
Ausgang des Wettlaufs entweder der erfolgreiche Abschluss oder der endgültige Fehlschlag eintritt.
Der Verlierer des Rennens darf dann keine Änderung mehr vornehmen, da der Gewinner den Job bereits in einen finalen
Zustand überführt hat.
Für \ref{zus:endzustaende} wird zusätzlich zur mehrfachen Ausführung protokolliert, wie viele Durchläufe auf den jeweiligen Ausgang entfallen,
um sicherzustellen, dass das Rennen tatsächlich beide möglichen Ausgänge erzeugt und nicht durch eine feste Reihenfolge vorentschieden ist.

\subsection{Recovery nach Worker-Ausfall (\ref{zus:verarbeitungsgarantie}, \ref{zus:idempotenz}, \ref{zus:terminierung})}
\label{subsec:recovery_fehlersimulation}
Zum Nachweis von \ref{zus:verarbeitungsgarantie} wird ein tatsächlicher Prozessabbruch herbeigeführt.
Dazu wird das Backend gestartet, aber die Anzahl laufender Worker auf 0 gesetzt.
Anschließend wird eine zweite Instanz des Systems als eigener Prozess gegen dieselbe Datenbank gestartet.
Dieser Prozess verfügt über einen laufenden Worker.
Aus dem Test heraus wird ein Job erzeugt und bis zu dessen Übernahme durch den Worker-Prozess abgewartet.
Sobald der Job in Bearbeitung ist, wird der Prozess hart beendet, ohne dass ein geordnetes Herunterfahren stattfindet.
Daraufhin wird verifiziert, dass der Worker-Prozess den Job tatsächlich nicht beenden konnte.
Anschließend wird ein neuer, unabhängiger Prozess mit laufendem Worker gestartet.
Der \texttt{RecoveryService} sollte den Job im nächsten Durchlauf als verwaist erkennen und erneut zur Bearbeitung ausstellen.
Der neue Worker sollte den verwaisten Job daraufhin übernehmen.
Zusätzlich zum Erreichen des Zustands \textit{Succeeded} wird für \ref{zus:idempotenz} geprüft, dass der gespeicherte Ergebnisinhalt dem
tatsächlich erfolgreichen, zweiten Verarbeitungsversuch entspricht und keine Reste des abgebrochenen ersten Versuchs enthält.
Funktioniert die Recovery wird außerdem gezeigt, dass ein Job nach Ablauf des Leases aus dem aktiven Zustand befreit wird,
sodass er dort nicht hängen bleibt (\ref{zus:terminierung}).

Der Test zu \ref{zus:idempotenz} (a) prüft, dass die Berechtigungen eines Workers nach der Recovery erlöschen.
Dazu wird nach Ablauf des Leases ein Recovery-Durchlauf aktiviert und gewartet, bis der Job durch einen anderen Worker
übernommen wird.
Daraufhin versucht der ursprüngliche Besitzer ein Ergebnis zu speichern.
Dieser Versuch wird auch nach Fertigstellung des Jobs erneut wiederholt.

Zusätzlich untersucht \ref{zus:idempotenz} (b), ob der Worker die Berechtigung nach Ablauf des Leases auch ohne vorherigen Recovery-Durchgang
verliert.
Zur Überprüfung wird ein Job mit Status \textit{Running} und abgelaufenem Lease im Job-Store persistiert.
Anschließend versucht der eingetragene Besitzer den Job über \textit{finishJob} zu beenden.

Mittels \texttt{@MockitoSpyBean} wird die Fachlogik eines einzelnen, gezielt ausgewählten Jobs über einen \texttt{JobProcessor}-Spy
so modifiziert, dass zum Test von \ref{zus:terminierung} jeder Verarbeitungsversuch mit einer Exception fehlschlägt.
Vor der eigentlichen Untersuchung wird zunächst ein Kontrolljob regulär verarbeitet, um sicherzustellen, dass der konfigurierte
Spy die grundsätzliche Funktionsfähigkeit des Workers nicht beeinträchtigt.
Anschließend wird geprüft, dass der fehlschlagende Job nach genau der konfigurierten Maximalanzahl der Wiederholungsversuche
in den Zustand \textit{Failed} überführt wird und kein weiterer Verarbeitungsversuch mehr stattfindet.
Da \ref{zus:terminierung} den Übergang in einen Endzustand nach endlicher Zeit fordert, wird für die maximale Zeit bis
zur endgültigen Beendigung eines dauerhaft fehlschlagenden Jobs eine obere Schranke definiert.
Diese ergibt sich aus der maximalen Anzahl der Verarbeitungsversuche sowie der für einen einzelnen Verarbeitungszyklus maximal
zu erwartenden Wartezeit:
\[
    T_{\mathrm{max}}
    =
    A_{\mathrm{max}}
    \cdot
    \left(
        T_{\mathrm{lease}}
        +
        T_{\mathrm{recovery}}
        +
        T_{\mathrm{idle}}
    \right)
\]
Dabei bezeichnet \(A_{\mathrm{max}}\) die maximal zulässige Anzahl von Verarbeitungsversuchen, bevor ein Job in \textit{Failed} überführt wird,
\(T_{\mathrm{lease}}\) die konfigurierte Lease-Dauer bis ein Job als verwaist gilt,
\(T_{\mathrm{recovery}}\) das Recovery-Intervall und
\(T_{\mathrm{idle}}\) das Warteintervall eines verfügbaren Workers zwischen zwei Übernahmeversuchen.
Um die Testlaufzeit zu optimieren, werden für ausschließlich diesen Test das Recovery-Intervall und die Lease-Dauer auf 1s reduziert.
In Kombination mit den in \cref{sec:versuchsaufbau} beschriebenen Parametern ergibt sich daraus für den Test folgende Konfiguration:
\begin{align*}
    A_{\mathrm{max}} &= 5\\
    T_{\mathrm{lease}} &= 1\,\mathrm{s},\\
    T_{\mathrm{recovery}} &= 1\,\mathrm{s},\\
    T_{\mathrm{idle}} &= 2\,\mathrm{s}.
\end{align*}
Daraus folgt, dass für die Zeit bis zum Übergang in einen finalen Status die obere Schranke von \(T_{\mathrm{max}} = 20\,s\) gilt.

\section{Ergebnisse}
\label{sec:ergebnisse}
Die im Folgenden dargestellten Beobachtungen stützen sich auf die tatsächlich protokollierten Testläufe.
Gekürzte, für die jeweilige Beobachtung repräsentative Auszüge der Logs finden sich im Anhang in \cref{sec:testlogs}.
Die vollständigen Protokolle aller Durchläufe sind unter \texttt{src/main/resources/test/Testlogs.md} in den Begleitdateien
beziehungsweise in dem im Anhang in \cref{sec:quellcode} verlinkten GitHub-Profil beigefügt.

\paragraph{Atomare Persistenz und Validität (\ref{zus:atomare_persistenz})}
\label{erg:atomare_persistenz}
Beide Requests wurden vom Backend mit dem HTTP-Status \texttt{400 Bad Request} abgelehnt, der erste aufgrund der
syntaktischen, der zweite aufgrund der fachlichen Validierung.
Die Anzahl der im Job-Store vorhandenen Jobs war vor und nach beiden Anfragen identisch, folglich kam es für keine der beiden
ungültigen Requests zu einer Persistierung des Jobs.
Damit wurde das für \ref{zus:atomare_persistenz} definierte Prüfkriterium erfüllt.
Die Atomarität der Persistierung wird durch die Verwendung von \texttt{@Transactional} und die \gls{acid}-Garantien von
PostgreSQL sichergestellt.
Das korrekte Rollback-Verhalten bei einem Fehler innerhalb der Transaktion wird bereits durch den Test von \ref{zus:atomare_wirkung} geprüft.
Da sich die transaktionale Umsetzung technisch nicht unterscheidet, wird hier auf einen separaten Test verzichtet.

\paragraph{Duplikaterkennung bei der Auftragserfassung (\ref{zus:duplikaterkennung})}
\label{erg:duplikaterkennung}
Im ersten Ausführungsversuch wurde der Job erfolgreich persistiert, bevor die simulierte \texttt{IOException} die Antwort abfing.
Der clientseitige \texttt{RetryExecutor} übertrug daraufhin dieselbe Anfrage samt identischem Idempotenzschlüssel erneut.
Das Backend erkannte dabei die Verletzung des entsprechenden \texttt{UNIQUE}-Constraints und behandelte den zweiten Aufruf
als Idempotenzfall, ohne diesen seinerseits erneut zu wiederholen.
Der clientseitige Retry wurde also tatsächlich ausgeführt, während der resultierende Datenbankfehler auf Backendseite
korrekt als nicht retrybar eingestuft wurde.
Der zweite Client-Aufruf lieferte den HTTP-Status \texttt{200 OK} und denselben, bereits zuvor erzeugten Job zurück.
Während des gesamten Tests erhöhte sich die Anzahl der Jobs im Job-Store lediglich von null auf eins.
Das für \ref{zus:duplikaterkennung} definierte Prüfkriterium ist damit erfüllt: Eine erneute Übermittlung derselben
Anfrage, wie hier beispielsweise nach einer verlorenen Antwort, führt nicht zu einer zweiten Persistierung, sondern liefert
den bereits vorhandenen Job zurück.

\paragraph{Exklusive Rechte (\ref{zus:exklusive_rechte})}
\label{erg:exklusive_rechte}
Bei der Untersuchung des Claimings (Testfall a) konnte in allen Durchläufen jeweils genau ein Worker den Job
erfolgreich übernehmen, erkennbar am einmaligen Zustandsübergang von \textit{Queued} nach \textit{Running}.
Die übrigen konkurrierenden Worker erhielten in keinem der Durchläufe einen Job.
Folglich wurde ein mehrfacher Übergang desselben Jobs nach \textit{Running} in keinem Durchlauf beobachtet, was bedeutet,
dass das Claiming exklusiv ablief.

Bei der Untersuchung der Ergebnisspeicherung (Testfall b) konnte in allen Durchläufen jeweils nur der aktuell
berechtigte Worker den Job erfolgreich abschließen.
Der konkurrierende Aufruf mit abweichender Worker-ID wurde je nach Scheduling der Anfragen entweder abgewiesen, da der
Job schon nicht mehr im Zustand \textit{Running} war, oder weil die Worker-ID nicht mit dem gespeicherten
\texttt{claimed\_by}-Wert übereinstimmte.
Nach jedem Testlauf existierte genau ein persistiertes \texttt{JobResult}, dessen Inhalt dem vom berechtigten Worker
übermittelten Ergebnis entspricht.
Damit wurde das für \ref{zus:exklusive_rechte} definierte Prüfkriterium in beiden Testfällen erfüllt: Ein
verfügbarer Job wird höchstens einem Worker zur Verarbeitung zugeordnet und von zwei konkurrierenden Abschlussversuchen
kann höchstens einer das Ergebnis persistieren.

\paragraph{Atomare Wirkung von Ergebnissen (\ref{zus:atomare_wirkung})}
\label{erg:atomare_wirkung}
Nach dem durch den Test-Hook erzwungenen Fehlschlagen des Aufrufs befand sich der Job weiterhin im Zustand
\textit{Running} und ein \texttt{JobResult} wurde nicht persistiert.
Die zuvor im Persistenzkontext vorgenommene Statusänderung zu \textit{Succeeded} wurde vollständig verworfen.
Das für \ref{zus:atomare_wirkung} definierte Prüfkriterium ist damit erfüllt: Tritt innerhalb der Transaktion, wie beispielsweise
zwischen Statusänderung und Ergebnispersistierung, ein Fehler auf, bleiben keine Teiländerungen in der Datenbank bestehen.
Die beobachteten Zustände bestätigen damit die durch die \gls{acid}-Garantien und das Transaktionsmodell umgesetzte atomare Wirkung.

\paragraph{Finalität von Endzuständen (\ref{zus:endzustaende})}
\label{erg:endzustaende}
In 47 der 100 Durchläufe gewann \textit{finishJob}.
Der Job erreichte den Zustand \textit{Succeeded}, die Ergebnisdaten wurden persistiert und der konkurrierende
\textit{recoverJob}-Aufruf wurde aufgrund des bereits erreichten Endzustands folgenlos beendet.
In den übrigen 53 Durchläufen gewann \textit{recoverJob}.
Da die maximale Anzahl an Verarbeitungsversuchen bereits erreicht war, wurde der Job in den Zustand \textit{Failed} überführt
und der konkurrierende \textit{finishJob}-Aufruf wurde abgewiesen.
Damit traten beide möglichen Ausgänge des konkurrierenden Zugriffs in den wiederholten Testläufen tatsächlich auf.
Anhand der protokollierten Zustandsübergänge jedes Durchlaufs wurde zusätzlich sichergestellt, dass in keinem der
100 Durchläufe ein Übergang zurück aus einem bereits erreichten Endzustand nach \textit{Queued} oder \textit{Running} erfolgte.
Die beobachteten Übergänge beschränkten sich durchgehend auf \textit{Queued} nach \textit{Running} sowie
anschließend entweder auf \textit{Running} nach \textit{Succeeded} oder \textit{Running} nach \textit{Failed}.

Die Finalität lässt sich zusätzlich konstruktiv argumentieren.
Im System existieren genau drei Stellen, an denen der Status eines Jobs verändert wird: \textit{claimNextJob}, \textit{finishJob}
und \textit{recoverJob}.
Für alle drei Stellen bestehen Vorbedingungen, die eine Änderung ausschließen, sobald ein Job bereits \textit{Succeeded}
oder \textit{Failed} erreicht hat. \textit{claimNextJob} berücksichtigt ausschließlich Jobs im Zustand \textit{Queued},
während \textit{finishJob} und \textit{recoverJob} ihre Verarbeitung abbrechen, wenn der Job nicht den Zustand \textit{Running} besitzt.
Da der relevante Datensatz bei allen drei Operationen mit \texttt{FOR UPDATE} innerhalb einer Transaktion gesperrt wird,
werden Lesen des Status, Prüfung der Vorbedingung und Zustandsänderung isoliert und atomar gegenüber konkurrierenden Zugriffen ausgeführt.
Eine zweite Transaktion mit altem Status kann die Prüfung daher nicht umgehen.
Da sämtliche statusverändernden Operationen durch diese Vorbedingungen eingeschränkt sind, folgt daraus, dass ein einmal
erreichter Endzustand nicht mehr verlassen werden kann.

\paragraph{Automatische Wiederholungsversuche (\ref{zus:retries})}
\label{erg:retries}
In Testfall (a) schlug der erste Zugriffsversuch aufgrund der nicht verfügbaren Datenbankverbindung fehl.
Der \texttt{RetryExecutor} klassifizierte den Fehler als retrybar und führte drei weitere Ausführungsversuche mit
Wartezeiten von 200\,ms, 400\,ms und 800\,ms durch.
Nach der Wiederherstellung der Datenbankverbindung schloss der vierte Versuch erfolgreich ab, sodass die Anfrage mit dem HTTP-Status
\texttt{201 Created} beantwortet wurde und der zurückgegebene Idempotenzschlüssel dem der ursprünglichen Request entsprach.
Ergänzend wurde als Gegenprobe geprüft, wie sich das System bei einem anhaltenden als transient klassifizierten Fehler verhält.
Die Anfrage wurde nach Erschöpfung aller konfigurierten Wiederholungsversuche mit dem HTTP-Status \texttt{500 Internal Server
Error} beantwortet und es wurde kein erneuter Retry gestartet.

Im zweiten Testfall führte die erneute Verwendung eines bereits vorhandenen Idempotenzschlüssels zu einer
\texttt{DataIntegrityViolationException}.
Der \texttt{RetryExecutor} stufte diese als nicht retrybar ein, sodass kein weiterer Ausführungsversuch erfolgte.
Der Controller behandelte den Fehler als Idempotenzfall und beantwortete die Anfrage mit \texttt{200 OK}.
Das für \ref{zus:retries} definierte Verhalten wurde damit tatsächlich beobachtet.

\paragraph{Worker-Ausfall und Wiederaufnahme (\ref{zus:verarbeitungsgarantie}, \ref{zus:idempotenz}, \ref{zus:terminierung})}
\label{erg:verarbeitungsgarantie}
Der gestartete Worker-Prozess übernahm den erzeugten Job und setzte dessen Status auf \textit{Running}.
Nach dem gezielten, harten Abbruch des Prozesses befand sich der Job weiterhin im Status \textit{Running}.
Der zunächst gestartete zweite Worker-Prozess übernahm den Job nicht sofort, sondern erst nachdem der Recovery-Service
den verwaisten Job beim nächsten Zyklus anhand des abgelaufenen Leases erkannte und ihn in den Zustand \textit{Queued} zurückführte.
Der zweite Worker übernahm den Job daraufhin und verarbeitete ihn erfolgreich.
Der Job erreichte den Status \textit{Succeeded} mit einem \texttt{attempt\_count} von zwei, was die ursprüngliche
Verarbeitung durch den ersten und die erneute Verarbeitung durch den zweiten Worker bestätigt.
Der gespeicherte Ergebnisinhalt entsprach dabei nachweislich dem tatsächlich erfolgreichen, zweiten Verarbeitungsversuch
und enthielt keine Reste des unterbrochenen ersten Versuchs.
Damit wurde das für \ref{zus:verarbeitungsgarantie} definierte Verhalten beobachtet: Der Ausfall eines Workers
führt nicht zum Verlust eines bereits persistierten Jobs, sondern zur automatischen Wiederfreigabe durch den Recovery-Service
und der anschließenden erneuten Verarbeitung.
Die Persistenz des Jobs über den Prozessabbruch hinweg wird dabei durch die Durability-Eigenschaft von PostgreSQL abgesichert.
Während die Recovery dafür sorgt, dass der Job nicht in der Verarbeitung hängen bleibt, sorgt \gls{acid} dafür, dass die
Jobdaten selbst nicht verloren gehen.
Dass der Job zudem mit inhaltlich korrektem, nicht durch den abgebrochenen ersten Versuch verfälschtem Ergebnis abschließt,
liefert zugleich einen Teilbeleg für \ref{zus:idempotenz}.
Der beobachtete Übergang aus \textit{Running} zurück nach \textit{Queued} nach Ablauf des Leases zeigt darüber hinaus einen
Teilaspekt von \ref{zus:terminierung}, nämlich dass ein Job nicht dauerhaft in einem aktiven Zustand verbleibt.

\paragraph{Eindeutigkeit des Ergebnisses bei abgelaufenem Lease (\ref{zus:idempotenz})}
\label{erg:idempotenz}
Im Testszenario (a) wurde der Job nach explizit ausgelöstem Recovery-Vorgang in den Zustand \textit{Queued} überführt
und von einem anderen Worker übernommen.
Der zuvor berechtigte, jetzt aber überholte Worker wurde bei seinem anschließenden \textit{finishJob}-Aufruf abgewiesen.
Der neue Worker schloss den Job dagegen erfolgreich ab.
Auch der \textit{finishJob}-Aufruf des überholten Workers nach der Fertigstellung wurde abgelehnt.
Insgesamt wurde genau ein \texttt{JobResult} für den Job persistiert.

In Testszenario (b) wurde der Aufruf von \textit{finishJob} durch den weiterhin eingetragenen Worker aufgrund des abgelaufenen
Leases mit einer \texttt{LeaseExpiredException} abgewiesen und kein Ergebnis persistiert.
Damit ist sowohl der Fall eines bereits durch Recovery übernommenen als auch der Fall eines abgelaufenen, aber noch nicht
durch die Recovery wiederhergestellten Jobs abgedeckt.
In beiden Fällen konnte ein verspäteter Worker kein Ergebnis mehr speichern, während der gültige Verarbeitungsversuch aus
Testfall (a) genau ein Ergebnis erzeugte.
Anknüpfend an die im vorangegangenen Abschnitt beobachtete inhaltliche Korrektheit des gespeicherten Ergebnisses nach
einem echten Prozessabbruch, zeigt dies den verbliebenen zu zeigenden Teil von \ref{zus:idempotenz}.
Das für \ref{zus:idempotenz} definierte Prüfkriterium wird damit erfüllt: Für einen Job wird höchstens ein Ergebnis
persistiert und ein verspäteter Worker wird an der Ergebnispersistierung gehindert.
Das beobachtete Verhalten realisiert damit die At-Most-Once-Eigenschaft für die Ergebnisspeicherung.
Zusammen mit der für \ref{zus:verarbeitungsgarantie} überprüften Wiederaufnahme nach einem Worker-Ausfall wird die
angestrebte At-Least-Once-Verarbeitung angenähert.

\paragraph{Terminierung in endlicher Zeit (\ref{zus:terminierung})}
\label{erg:terminierung}
Der zunächst verarbeitete Kontrolljob erreichte erwartungsgemäß den Status \textit{Succeeded} und bestätigte damit,
dass der konfigurierte \texttt{JobProcessor}-Spy die grundsätzliche Funktionsfähigkeit des Workers nicht beeinträchtigt.
Für den anschließend erzeugten Testjob trat bei jedem Übernahmeversuch der simulierte Fachfehler auf.
Nach jedem fehlgeschlagenen Versuch lief das Lease ab, woraufhin der Recovery-Service den Job erneut in den Zustand \textit{Queued} überführte.
Dieser Ablauf wiederholte sich, bis die konfigurierte maximale Anzahl an Verarbeitungsversuchen erreicht war.
Nach dem letzten fehlgeschlagenen Versuch überführte der Recovery-Service den Job in den Endzustand \textit{Failed}.
Der abschließende \texttt{attempt\_count} entsprach dabei genau dem konfigurierten Wert \texttt{maxAttempts}.
Das für \ref{zus:terminierung} definierte Kriterium ist damit erfüllt: Der dauerhaft fehlschlagende Job verbleibt
nicht unbegrenzt im Wechsel zwischen \textit{Running} und \textit{Queued}, sondern wird nach Erreichen der
maximalen Anzahl von Verarbeitungsversuchen endgültig in den Zustand \textit{Failed} überführt.
Die hierfür benötigte Zeit betrug 12,179\,s und lag damit unterhalb der aus den konfigurierten Zeitparametern
abgeleiteten oberen Schranke von 20\,s.

Für die Terminierung ist neben der Begrenzung der Anzahl der Verarbeitungsversuche auch die Reihenfolge der Jobauswahl relevant.
Die Auswahl erfolgt wie in \cref{subsec:impl_persistenz} beschrieben anhand des Erstellungszeitpunkts durch die Sortierung
der Jobs innerhalb der Datenbankanfrage nach dem Attribut \texttt{created\_at}.
Dadurch kann es nicht passieren, dass ein Job dauerhaft durch später erstellte Jobs verdrängt wird.

\section{Bewertung der Hypothesen}
\label{sec:hypothesenbewertung}
\textbf{\ref{item:h1}} nimmt an, dass sich durch die Kombination eines persistenten Statusmodells mit
transaktionsbasiertem Zugriff und Locking sowie Idempotenz-Mechanismen ein System entwerfen lässt, das auch bei
unerwarteten Prozessabbrüchen und Teilausfällen keine Jobs verliert und keine falschen beziehungsweise doppelten Ergebnisse
erzeugt.

Die durchgeführten Testfälle stützen diese Hypothese.
Die exklusive Beanspruchung eines Jobs sowie die exklusive Ergebnisspeicherung wurden für \ref{zus:exklusive_rechte} in
beiden Testfällen über jeweils 100 Wiederholungen ohne eine einzige Verletzung nachgewiesen.
Die für \ref{zus:atomare_wirkung} und \ref{zus:endzustaende} beobachteten Ergebnisse bestätigen zusätzlich, dass weder
Teiländerungen noch ein Rückfall aus einem bereits erreichten Endzustand auftreten.
Letzteres wird durch die konstruktive Betrachtung der drei zustandsverändernden Codepfade zusätzlich bestärkt.
Der tatsächliche, harte Absturz eines Worker-Prozesses führte ebenfalls nicht zum Verlust des betroffenen Jobs oder einem falschen Ergebnis.
Dass dabei zu keinem Zeitpunkt mehr als ein Ergebnis pro Job persistiert wurde, bestätigt zusätzlich \ref{zus:idempotenz}
und somit insgesamt \ref{item:h1}.

Einschränkend ist jedoch festzuhalten, dass nicht garantiert wird, dass für jeden persistierten Job ein erfolgreiches Ergebnis erzeugt wird.
Tritt bei der Verarbeitung wiederholt ein Fehler auf, erreicht der Job im Rahmen von \ref{zus:terminierung} nach endlich
vielen Versuchen bewusst einen definierten Fehlerzustand, statt unbegrenzt oft wiederholt zu werden.
Diese bewusste Grenze wird in \cref{sec:einschraenkungen} nochmals vertiefend eingeordnet.

\textbf{\ref{item:h2}} hat als Annahme, dass Retry-Strategien mit exponentiellem Backoff in Verbindung mit einer
Fehlerklassifikation in transiente und permanente Fehler die Robustheit des Systems gegenüber
kurzzeitigen Infrastrukturausfällen erhöhen, ohne dabei die definierten Zusicherungen zu verletzen.

Die Untersuchung von \ref{zus:retries} stützt diese Hypothese direkt.
Ein kurzzeitiger Datenbankausfall wurde korrekt als retrybar klassifiziert und mit exponentiell steigenden Wartezeiten
automatisch überbrückt, sodass die ursprüngliche Anfrage letztlich erfolgreich abgeschlossen wurde, ohne dass ein
manuelles Eingreifen erforderlich wurde.
Ein permanenter Fehler wurde dagegen korrekt als nicht retrybar erkannt und ohne einen einzigen weiteren Ausführungsversuch
sofort beantwortet.
Die Fehlerklassifikation griff also in beide Richtungen zuverlässig.
Die ergänzend durchgeführte Gegenprobe bezüglich \ref{zus:duplikaterkennung} und \ref{zus:retries} zeigt außerdem,
dass die Resistenzsteigerung nicht auf Kosten anderer Zusicherungen geht.
Ein durch einen Retry ausgelöster doppelter Persistierungsauftrag eines neuen Jobs wurde erfolgreich abgewiesen und eine
Anfrage mit einem als transient klassifizierten Fehler wurde nach Erschöpfung der konfigurierten Wiederholungsversuche
kontrolliert mit einem Fehlerstatus beantwortet, anstatt die Ausführung unbegrenzt weiter zu versuchen.

\textbf{\ref{item:h3}} behauptet, dass sich die definierten Zusicherungen durch automatisierte Tests mittels gezielter
Fehlerinjektion reproduzierbar überprüfen und nachweisen lassen.

Für jede der neun Zusicherungen Z-1 bis Z-9 wurde, wie in \cref{ch:methodik} methodisch festgelegt, mindestens
ein automatisierter Testfall mit gezielter Fehlerinjektion umgesetzt (vgl.\cref{tab:testmatrix}) und erfolgreich durchgeführt.
Sämtliche Testfälle laufen dabei gegen einen für jeden Durchlauf neu hergestellten, isolierten Ausgangszustand
und lassen sich beliebig oft wiederholt ausführen, ohne dass eine manuelle Anpassung erforderlich ist.
Besondere Relevanz für die Reproduzierbarkeit hat dabei die Wiederholung der Testfälle zu \ref{zus:exklusive_rechte}
und \ref{zus:endzustaende}.
Da diese im Gegensatz zu den übrigen Testfällen aufgrund der Abhängigkeit vom Thread-Scheduling nicht-deterministisch sind,
würde ein einmalig bestandener Testlauf nur eines von potenziell mehreren möglichen Ergebnissen zeigen.
Die Wiederholung bestätigt hier, dass die Kriterien unabhängig von der konkreten Schedule zuverlässig eingehalten werden,
statt nur einmalig zufällig zuzutreffen.
Dabei wurde auch nachgewiesen, dass tatsächlich unterschiedliche Schedules auftraten.
Ein Fehler, der nur unter einer bestimmten, seltener auftretenden Ausführungsreihenfolge sichtbar würde, wäre dadurch mit
hoher Wahrscheinlichkeit erkannt worden.
\Ref{item:h3} wird somit als erfüllt betrachtet.

Insgesamt stützen die in \cref{sec:ergebnisse} dokumentierten Ergebnisse alle drei Hypothesen.
Kein durchgeführter Testfall lieferte einen Befund, der einer der drei Hypothesen widerspricht.
Eine Diskussion des damit verbundenen Erkenntnisgewinns sowie dessen Grenzen erfolgt im folgenden Kapitel.
